\section{Prometheus Architecture}
\label{sec:architecture-main}

Prometheus investigates a trainable computational Block whose internal
operations, calls, and memory are learned together from task feedback. The
current development architecture uses a continuous state, a bank of nonlinear
residual primitives, and content-addressed memory. This revises the earlier
categorical and loop-based implementations without transferring their
protection guarantees to the new model. Its ingredients are established neural
mechanisms; the question is whether their particular organization supports
useful acquisition, composition, and retention at competitive total cost.

\subsection{State, inputs, and executable primitives}

Let $h_{t,s}\in\mathbb R^H$ denote the state before call slot $s$ of input $t$.
Memory contains keys $\mathsf K_t\in\mathbb R^{M\times Q}$ and values
$\mathsf V_t\in\mathbb R^{M\times H}$. An observed initial-state identifier
selects a learned embedding $h_{1,0}=E_\theta[x]$. Keys begin at a learned,
task-independent matrix and values at zero for each example. There is no
persistent memory across examples in this implementation.

The raw input feature $u_t\in\mathbb R^T$ is a token embedding for symbolic
inputs or the output of a learned modality adapter. A learned slot embedding
gives $x_{t,s}=u_t+e_s$. All slots reuse the same primitive bank and caller.
During a token, memory is fixed while the state changes. Each call computes
\begin{align}
 r_{t,s}&=\operatorname{softmax}\!\left(
   q_\theta(h_{t,s},x_{t,s})\mathsf K_t^\top/\sqrt Q\right)\mathsf V_t,
   \\[-1mm]
 \alpha_{t,s}&=\operatorname{softmax}
   f_\theta(h_{t,s},x_{t,s},r_{t,s}),\\
 z_{t,s}&=[h_{t,s};x_{t,s};r_{t,s}],\\
 \delta_{t,s,k}&=U_k\operatorname{GELU}(D_k z_{t,s}+b_k),\\
 h_{t,s+1}&=\operatorname{LayerNorm}\!\left(
    h_{t,s}+\sum_{k=1}^{K}\alpha_{t,s,k}\delta_{t,s,k}\right).
 \label{eq:main-block-computation}
\end{align}
Here $D_k\in\mathbb R^{R\times(2H+T)}$ and
$U_k\in\mathbb R^{H\times R}$. Each primitive is an executable nonlinear
state update with its own learned parameters. The bank has no supplied
task-specific operation or explicit identity choice. A residual path replaces
repeated row-stochastic state propagation; neither normalization nor low rank
implies invertibility or exact preservation of information.

\subsection{Content memory and composition}

After the last call, write features concatenate the final state, raw input,
and last read vector: $\bar z_t=[h_{t,S};u_t;r_{t,S-1}]$. Learned projections
produce an address, gate, candidate key and candidate value:
\begin{align}
 a_t&=\operatorname{softmax}
    (w_\theta(\bar z_t)\mathsf K_t^\top/\sqrt Q),
       &g_t&=\sigma(g_\theta(\bar z_t)),\\
 \widehat k_t&=\tanh(k_\theta(\bar z_t)),
       &\widehat v_t&=v_\theta(\bar z_t).
\end{align}
With $\beta_{tj}=g_ta_{tj}$, each slot updates as
\begin{equation}
 (\mathsf K_{t+1,j},\mathsf V_{t+1,j})=
 (1-\beta_{tj})(\mathsf K_{t,j},\mathsf V_{t,j})+
 \beta_{tj}(\widehat k_t,\widehat v_t).
\end{equation}
The next token starts from $h_{t+1,0}=h_{t,S}$. Thus later calls depend on
the actual intermediate computation and retained content. Order-sensitive
composition follows from executing successive nonlinear updates; no
commutativity constraint is imposed. The terminal distribution is
$\operatorname{softmax}(W_Rh_{L,S}+b_R)$.

The development configuration has $H=64$, $T=Q=16$, $M=8$, $K=4$, $R=8$ and
one call per token. All dimensions are explicit resources, rather than counts
inferred from a hidden task generator. A static-caller control zeros only the
state and memory-read inputs to $f_\theta$; its primitives, memory, shapes and
parameter count are unchanged. It therefore retains contextual computation
and is stronger than a memoryless ablation.

\subsection{A common input interface}

Symbolic lookup and externally embedded inputs enter the same execution
function. Audio features or image patches require a learned projection to
$\mathbb R^T$, whose parameters and work are charged separately. No modality
adapter supplies operations, calls, intermediate targets, or task identities.
This common interface permits testing the same core across tasks; it does not
by itself establish transfer or continual retention. The complete Primary
architecture remains unfrozen, and present measurements are separately
identified development evidence.
